% ============================================================================
% Appendix E — Reproducibility recipe
% The data preparation pipeline, snapshot recovery sequence, step-by-step
% run order with gate verdicts, framework version pins, hardware envelope,
% and a verification checklist a reproducer uses to confirm the headline
% numbers match within tolerance.
% ============================================================================

\chapter{Reproducibility recipe}
\label{app:reproducibility}

\section{Data preparation}
\label{app:repro-data}

The five benchmark loaders that feed the realised evaluation panel are registered in \verb|eval/benchmarks.py| against the Hugging Face dataset hub identifiers. FEVER reads through \verb|lucadiliello/fever| under the development and training splits; TriviaQA reads through \verb|trivia_qa| under the reading-comprehension-no-context configuration; CommonsenseQA reads through \verb|commonsense_qa|; TruthfulQA reads through \verb|truthful_qa| under the generation configuration; StrategyQA reads through \verb|ChilleD/StrategyQA| with documented fallbacks to \verb|wics/strategy-qa| and the Eladsegal train-JSON mirror. Earlier loader registrations for Natural Questions and HotpotQA are preserved in the module for reproducibility of the cycle-zero diagnostics but do not feed the realised five-benchmark panel.

The six-way disjoint partition is constructed once at pool-builder time by \codepath{caem/benchmark\_splits.py::build\_all\_benchmark\_pools}, which is invoked by the experiment runner and writes the per-benchmark identifier manifest to \verb|outputs/full_run/dataset_splits.json|. Each benchmark contributes a $1{,}000$-sample seed pool for cold-start memory population, a $500$-sample calibration fold, a $500$-sample purity-validation slice, ten per-cycle stream-chunk slices (FEVER and TriviaQA at $2{,}000$ samples per cycle, CommonsenseQA at $700$ samples per cycle to fit the smaller available train pool of about $10{,}000$ samples after deduplication), a $300$-sample evaluation fold, and a test fold of up to $500$ samples (clamped to the available remainder where the source split admits fewer than the registered budget). The shuffle is seeded at forty-two on the training side and forty-three on the evaluation side, so the partitions are stable across re-invocations on the same source dataset.

Within-benchmark deduplication runs through the \verb|_dedupe_by_content_id| helper, which keys on the first $500$ characters of the question text under a sixteen-character SHA-256 hash. Cross-benchmark disjointness is verified by the \codepath{assert\_cross\_benchmark\_disjoint} assertion. Both checks run at pool-builder time; the runner re-stamps the same checks at every successful cycle boundary through \codepath{\_dedup\_inplace} and \codepath{assert\_disjoint\_calibration}, so any drift in the partition would be caught at the cycle close rather than at the metric-computation step.

The realised evaluation fold reads $300$ samples per benchmark across the five-benchmark panel, for a total of $1{,}500$ evaluation questions per cycle. The StrategyQA test fold has $387$ samples rather than the registered $500$ because the source split (the ChilleD StrategyQA test split) has only $687$ labelled samples, of which $300$ form the evaluation fold.

\section{Snapshot recovery}
\label{app:repro-snapshot}

The pre-main snapshot is the state of the run just before the main cycle loop. It is held in the private Hugging Face dataset repository \verb|aksaN000/caem-passage-index-21m| (access on request), which also holds the $64$-gigabyte Wikipedia passage index at its root, with a copy in the Drive folder's \texttt{9\_archive/pre\_main\_snapshot/}. Each bundle sits under \verb|pre_main_snapshot/<git-commit>/|: the anchor before the main run is \verb|c5b8066|, and the anchor after the retrieval-utility classifier was added is \verb|700b9e3|. A manifest in each bundle records the upload time, the source machine, the commit and a recovery hint.

The three-command recovery sequence runs from a fresh clone with the Hugging Face command-line interface authenticated to the private repository:
\begin{verbatim}
# 1. Fetch the index and the pre-main snapshot bundles
hf download aksaN000/caem-passage-index-21m \
    --include 'passages.*' 'pre_main_snapshot/*'

# 2. Stage the chosen anchor into outputs/
cp -r pre_main_snapshot/c5b8066/* outputs/

# 3. Pin the code to the matching commit
git checkout c5b8066    # or 700b9e3 for the RUC-era resume point
\end{verbatim}
After the third command the working tree holds the cold-start memory store, the cycle-zero memory and deferred-buffer snapshots, the cycle-zero calibration fold (per-sample nine-signal vectors and gold labels), the cycle-zero composite calibration JSON, the calibrated configuration JSON, the cycle-zero validation-gate verdict, the cycle-zero evaluation fold under the pristine generator, the post-calibration rescored evaluation fold, the dataset-splits manifest, and the run logs. The four artefacts that must be regenerated rather than pulled from the snapshot are the Qwen-2.5-3B-Instruct base weights (recovered from the Hugging Face model hub at deployment time), the per-cycle LoRA adapters at cycles one through five, the per-cycle self-improvement training artefacts, and the post-cycle-zero composite refits.

\section{Step-by-step run sequence}
\label{app:repro-steps}

The commands below reproduce the run from a fresh clone. What each stage does, and the gate it must pass, is described once in the stage graph of \Cref{app:runbook}; the run itself was driven by the top-level script that is now archived (\verb|archive/scripts_superseded/old_runners/run_phase1a.sh|), whose stages call these commands.

\paragraph{Pre-run setup:}
\begin{verbatim}
python scripts/prompt_smoke_test.py \
    --benchmarks fever triviaqa commonsense_qa
python scripts/build_alias_dict_from_benchmarks.py
python -m scripts.seed_cold_start
python -m scripts.run_experiment --output_dir outputs/cycle_0 \
    --num_cycles 0 --n_questions 500 --n_eval_questions 500 \
    --benchmarks fever triviaqa commonsense_qa truthfulqa strategyqa \
    --passage_index data/passage_index \
    --cold_start_memory outputs/cold_start_memory/memory_store
python scripts/fit_composite_calibration.py --shrinkage_alpha 0.6 \
    --cherian_boost --boost_C 0.01
python scripts/validate_composite_weights.py
\end{verbatim}
The retention-probe baseline and the snapshot upload are the stages \verb|step_19_2_slice| and \verb|step_gdrive_upload_pre_main| of the archived driver.

\paragraph{Main cycle loop:} A single invocation drives all six per-cycle sub-steps and resumes from the highest-cycle artefact set on disk:
\begin{verbatim}
python -m scripts.run_experiment \
    --output_dir outputs/full_run \
    --num_cycles 10 \
    --n_questions 5000 \
    --n_eval_questions 300 \
    --benchmarks fever triviaqa commonsense_qa truthfulqa strategyqa \
    --passage_index data/passage_index \
    --cold_start_memory outputs/cold_start_memory/memory_store \
    --verifier_backend minicheck \
    --resume_from_cycle <N>
\end{verbatim}
The six steps of each cycle and the rollback contract are described in \Cref{app:runbook}; the time per cycle is given in \Cref{app:repro-hardware}.

\paragraph{Post-run analysis:} The baselines and every reported number are produced by two commands.
\begin{description}[style=nextline,leftmargin=0pt]
\item[Baseline panel under the comparison protocol:] The baseline re-run of \Cref{sec:setup-protocol} on one rented $32$-gigabyte graphics processor: retrieval of the evaluation passages by exact search, B1, B2, B5 and B8 on all five benchmarks, B3 and B4 on TruthfulQA, FLARE, the evaluation of the vanilla fine-tune, and the verifier pass over every baseline answer with the cycle-three adapter and composite. Each stage can be run on its own.
\begin{verbatim}
bash scripts/rental/final_version.sh all
\end{verbatim}
\item[Final numbers, tables and figures:] On CPU in about two minutes: writes the final correctness label into every answer file, recomputes CHM in the model and delivered views and the status-tag accuracy, runs the paired tests, and writes every chapter-five table and figure and the result tables of the appendices. The option \texttt{--judge} first re-judges every TruthfulQA answer with the language-model judge (requires \texttt{ANTHROPIC\_API\_KEY}); the option \texttt{--check} fails if any committed output changes.
\begin{verbatim}
bash scripts/make_final_results.sh --check
\end{verbatim}
\end{description}

\section{Version pinning}
\label{app:repro-versions}

The environment is pinned in \verb|requirements.txt|, with PyTorch installed first from the CUDA wheel index; the exact versions of the baseline re-run environment are frozen in \verb|requirements-lock-rental-2026-10-18.txt|. The main run used the versions below.

\begin{table}[!ht]
\centering\footnotesize
\setlength{\tabcolsep}{5pt}
\renewcommand{\arraystretch}{1.1}
\rowcolors{2}{gray!8}{white}
\begin{tabular}{@{}p{5.0cm}p{3.0cm}p{5.5cm}@{}}
\toprule
\textbf{Package} & \textbf{Version} & \textbf{Role} \\
\midrule
Python                & $3.12$                   & interpreter \\
CUDA toolkit          & $13.0$ (driver $13.1$)   & RTX 5090 Blackwell compute capability \texttt{sm\_120} \\
\midrule
torch                 & $2.10.0\text{+cu130}$    & tensor backend, bf16 forward path, MC-dropout sampler \\
transformers          & $4.57.6$                 & Qwen-2.5-3B-Instruct loader, generate, MiniCheck encoder \\
peft                  & $0.19+$                  & LoRA wrapper for the self-improvement fine-tune \\
accelerate            & $0.34+$                  & device placement during the fine-tune step \\
\midrule
datasets              & $3.0+$                   & benchmark loaders and retention-probe shuffles \\
faiss-cpu             & $1.8+$                   & passage index and episodic memory store \\
sentence-transformers & $4.1.0$                  & SBERT memory encoder and cross-encoder reranker \\
\midrule
scipy                 & $1.11+$                  & Wilcoxon test, temperature fit (L-BFGS) \\
scikit-learn          & $1.4+$                   & per-signal isotonic regression and boost-layer logistic \\
\midrule
bitsandbytes          & $0.43+$                  & eight-bit AdamW optimiser path \\
anthropic             & $0.108$ to below $1$     & TruthfulQA judge client \\
numpy, pandas, tqdm   & latest stable            & arithmetic, table aggregation, progress bars \\
\midrule
biber + biblatex      & TeX Live 2024+           & thesis bibliography compile \\
\bottomrule
\end{tabular}
\caption[Framework version pinning]{Framework version pinning. Lower bounds follow the requirements file; the lock file gives the exact versions of the re-run.}
\label{tab:app-repro-versions}
\end{table}

Seeding is partial, and the numbers are reproducible within the tolerances of \Cref{app:repro-verify} rather than bit for bit. The seed $42$ fixes the benchmark partitions (\Cref{app:repro-data}) and the order of each fine-tune; generation is greedy, so answers repeat given the same weights and inputs. The signals computed from sampled text (dropout agreement, chain similarity and chain entailment) are random draws by design, PyTorch's deterministic-algorithm mode is left off because the attention kernels are not bitwise deterministic, and TensorFloat-32 matrix multiplication is left on.

The two required environment variables are \verb|ANTHROPIC_API_KEY| for the TruthfulQA judge and \verb|HF_TOKEN| (equivalently \verb|HUGGINGFACE_HUB_TOKEN|) for the pull from the private pre-main snapshot repository. Optional environment variables \verb|HUGGINGFACE_HUB_CACHE|, \verb|CAEM_DEVICE|, and \verb|CUDA_VISIBLE_DEVICES| override the cache location, the compute device, and the GPU index respectively.

\section{Hardware envelope}
\label{app:repro-hardware}

\paragraph{Production envelope:} the run used a rented RTX 5090 with $32$ gigabytes of graphics memory, about $150$ gigabytes of disk and about $90$ gigabytes of host memory, at about $0.70$ US dollars per hour. Cycles zero to five of the main run took $6.0$, $29.0$, $35.4$, $43.2$, $24.3$ and $54.6$ hours according to the run log, about $193$ hours in all (\Cref{sec:env-experiment}); the growth comes mostly from the retroactive re-verification pass, which grows with the memory (about $38$ minutes for $431$ entries at cycle one, about $16$ hours for $10{,}214$ entries at cycle five). The verifier, which re-scores every answer, takes most of the time; the fine-tune itself is a small share. With the calibration phase, the earlier rejected runs and the analysis, the programme used several times the compute of the main run (\Cref{sec:env-experiment}). The baseline re-run of \Cref{sec:setup-protocol} ran later on a second rented machine of the same class.

\paragraph{Writeup envelope:} the thesis was written on a Windows Subsystem for Linux 2 environment with an RTX 3060 ($12$ gigabytes of graphics memory). The $64$-gigabyte passage index does not fit in its memory; all retrieval-side numbers in the thesis are read from cached evaluation JSON files under \verb|outputs/full_run/| rather than recomputed locally. The local writeup envelope hosts the thesis source, the table and figure compilation, and the signal-level audits over saved verifier traces, but does not host the live retrieval substrate.

\section{Verification of reproduced numbers}
\label{app:repro-verify}

A reproducer who has executed the full trajectory uses the ten-item checklist below to confirm that the headline numbers match the thesis within the stated tolerance. The per-benchmark exact match on $300$ questions carries a stochastic floor of about $\pm 0.5$ percentage points from the non-deterministic attention path; the pooled exact match over the five benchmarks tightens to about $\pm 0.3$ points; the composite hallucination metric and the storage precision tolerate about $\pm 1.0$ point; integer counts (commit pattern, win and defeat tallies) must match exactly when the results command is run on the same answer files. Items one to three, eight and ten are rebuilt from the saved answer files by the results command above, which reproduces them exactly.

\begin{description}[style=nextline,leftmargin=0pt]
\item[Pooled CAEM exact match:]
\begin{itemize}
\item Receipt: the final correctness label of every answer in \codepath{outputs/full\_run/eval/}, pooled over the five benchmarks (\Cref{tab:per-bench-trajectory}).
\item Expected: $0.517$ at C0, $0.535$ at C2, $0.508$ at C5; B1 $0.511$.
\item Tolerance: $\pm 0.3$ percentage points.
\end{itemize}
\item[Pooled composite hallucination metric:]
\begin{itemize}
\item Receipt: \codepath{outputs/analysis/chm\_judged.json}, model view (mean of the seven active failure flags).
\item Expected: $0.155$ at C0, $0.127$ at C2, $0.134$ at C5; B1 $0.132$.
\item Tolerance: $\pm 1.0$ percentage point.
\end{itemize}
\item[Per-benchmark exact match and CHM by cycle:]
\begin{itemize}
\item Receipt: the same record, per benchmark.
\item Expected: each cell of \Cref{tab:per-bench-trajectory}.
\item Tolerance: $\pm 0.5$ percentage points.
\end{itemize}
\item[Calibration-fold pooled storage precision band:]
\begin{itemize}
\item Receipt: \codepath{outputs/full\_run/pi\_store\_trajectory.csv}, column \texttt{cal\_fold\_pi\_store} across the committed-cycle rows.
\item Expected: band $[0.732, 0.746]$.
\item Tolerance: $\pm 1.0$ percentage point on each endpoint.
\end{itemize}
\item[Cycle-zero validation-gate composite Cohen's $d$:]
\begin{itemize}
\item Receipt: \codepath{outputs/analysis/validation\_gate\_locked.json}, field \texttt{d\_pooled\_train} (the locked composite on the validation fold, \Cref{sec:adj-gate}).
\item Expected: $0.717$.
\item Tolerance: $\pm 0.02$.
\end{itemize}
\item[Memory pool size at cycle five close:]
\begin{itemize}
\item Receipt: \codepath{outputs/full\_run/memory\_store\_cycle\_5.meta}, entry count.
\item Expected: $12\,175$ entries.
\item Tolerance: $\pm 5$ for content-hash collision drift.
\end{itemize}
\item[Realised commit pattern and abort cycle:]
\begin{itemize}
\item Receipt: \codepath{outputs/full\_run/stochastic\_equilibrium\_trajectory.csv}, columns \texttt{verdict} and \texttt{committed}.
\item Expected: commit pattern $(1, 1, 1, 0, 1)$ across C1 to C5, with cycle four aborted by the TriviaQA-test retention probe.
\item Tolerance: exact match required.
\end{itemize}
\item[Retrieval-utility classifier ablation:]
\begin{itemize}
\item Receipt: the classifier-off answers at cycle three in \codepath{outputs/phase1d\_eval/} against the main run at cycle three (\Cref{tab:ruc-ablation}).
\item Expected: pooled exact match $0.467$ off and $0.516$ on, a lift of $4.9$ percentage points.
\item Tolerance: $\pm 0.5$ percentage points.
\end{itemize}
\item[Status-tag accuracy:]
\begin{itemize}
\item Receipt: \codepath{outputs/analysis/chm\_judged.json}, field \texttt{tags} of each CAEM cycle (\Cref{tab:status-tags}).
\item Expected: verified answers correct $0.746$, $0.655$ and $0.667$ of the time at C0, C2 and C5.
\item Tolerance: $\pm 1.0$ percentage point.
\end{itemize}
\item[Significance win and defeat tallies:]
\begin{itemize}
\item Receipt: \codepath{outputs/analysis/paired\_baseline\_tests\_final.json} (\Cref{tab:sig-summary}).
\item Expected: exact match $9$ wins and $2$ defeats at C0, $9$ and $3$ at C2, $5$ and $3$ at C5; CHM $6$ and $17$ at C0, $13$ and $5$ at C2, $10$ and $4$ at C5, each out of $40$ contrasts.
\item Tolerance: exact match required.
\end{itemize}
\end{description}

A reproducer whose readings clear every item within the stated tolerance has matched the thesis to within the determinism floor of the compute envelope. An item that fails by a larger margin points to a determinism gap the runner did not absorb (the seed-everything helper is the first place to inspect), a version mismatch against \Cref{tab:app-repro-versions}, a hardware change not registered in the runner's overrides, or a genuine correction that the empirical chapter should adopt.
